\documentclass[11pt]{article}
\usepackage[a4paper,margin=22mm]{geometry}
\usepackage{fontspec}
\setmainfont{DejaVu Serif}
\setsansfont{DejaVu Sans}
\usepackage{microtype}
\usepackage{xcolor}
\usepackage{hyperref}
\usepackage{enumitem}
\usepackage{parskip}
\definecolor{Accent}{HTML}{971D32}
\definecolor{Muted}{HTML}{666666}
\hypersetup{colorlinks=true,urlcolor=Accent,linkcolor=Accent}
\setlist{leftmargin=1.6em,itemsep=0.3em,topsep=0.35em}
\setcounter{tocdepth}{2}
\renewcommand{\contentsname}{Contents}
\newcommand{\sectionrule}{\par\vspace{-0.5em}\noindent\textcolor{Accent}{\rule{\linewidth}{0.6pt}}\par}
\begin{document}

\begin{center}
{\sffamily\bfseries\color{Accent} TOEFL WORD OF THE DAY\par}
\vspace{8mm}
{\fontsize{42}{48}\selectfont\bfseries devious\par}
\vspace{2mm}
{\Large\sffamily\color{Accent} /ˈdiːviəs/\par}
\vspace{2mm}
{\small\sffamily Local audio phrase: devious\par}
{\small\sffamily Audio file: audios/devious.mp3\par}
\vspace{2mm}
{\large\sffamily\color{Muted} adjective\par}
\vspace{5mm}
{\sffamily dishonest and cleverly indirect \textbullet\ indirect or winding\par}
\vspace{6mm}
{\large Devious usually describes dishonest or manipulative behavior, but it can also describe a route that is indirect or winding.\par}
\vspace{6mm}
\begin{minipage}{0.84\textwidth}
\itshape “Investigators uncovered a devious scheme to avoid paying taxes.”
\end{minipage}\par
\vspace{10mm}
\textbf{Date:} August 25th 2026\quad\textbullet\quad \textbf{Level:} B1--mid-B2 TOEFL
\vfill
Created and compiled by \textbf{Amir Kooshky}\par
\vspace{2mm}
\href{https://t.me/KooshkyTOEFL}{Telegram Channel}\quad\textbullet\quad
\href{https://www.instagram.com/kooshkytoefl}{Instagram}\quad\textbullet\quad
\href{https://t.me/KooshkyTOEFL\_pv}{Personal Telegram}
\end{center}
\newpage
\tableofcontents
\newpage

\section{Meanings and usage}
\sectionrule
\subsection{Meaning 1: Dishonest and cleverly indirect}
\textbf{Part of speech:} adjective

\textbf{IPA:} /ˈdiːviəs/

\textbf{Pronunciation phrase:} devious

\textbf{Local audio file:} audios/devious.mp3

\textbf{Register:} neutral to formal; strongly negative and common in descriptions of dishonest, manipulative, or secretive behavior

\textbf{Definition:} using dishonest, secretive, or clever methods to get what you want

\textbf{In simpler words:} sneaky and dishonest

\textbf{Typical pattern:} devious + tactic, scheme, method, plan, motive, or person

\subsubsection{Natural examples}
\begin{enumerate}
  \item The manager used devious tactics to hide the true cost of the project.
  \item Investigators uncovered a devious scheme to avoid paying taxes.
  \item The character is intelligent but devious and frequently manipulates people for his own benefit.
  \item The company was accused of using devious methods to discourage customers from canceling their subscriptions.
  \item His devious plan depended on keeping several important facts secret.
  \item The politician's opponents described the campaign strategy as devious and misleading.
  \item She suspected that he had a devious motive for offering to help.
  \item A devious negotiator may appear cooperative while secretly trying to gain an unfair advantage.
\end{enumerate}

\subsubsection{Nuance and implication}
\textbf{Central nuance:} Devious suggests dishonesty combined with cleverness or indirectness rather than simple open wrongdoing.

\textbf{Usual implication:} A devious person often hides their true intentions and uses complicated or manipulative methods to gain an advantage.

\textbf{Compared with straightforward:} A straightforward person or method is open, direct, and easy to understand. A devious one achieves its purpose through hidden, indirect, or dishonest means.

\subsubsection{Grammar notes}
\textbf{How to use it:} Use it before a noun, as in a devious plan, or after a linking verb, as in His methods were devious.

\textbf{What can follow it:} It commonly describes people, tactics, schemes, plans, methods, strategies, and motives.

\textbf{Position in a sentence:} Words such as extremely, highly, particularly, and remarkably can appear before devious.

\subsubsection{Useful patterns}
\textbf{Pattern:} a devious plan or scheme

\textbf{Use:} Use this for a dishonest plan involving secrecy or manipulation.

\textbf{Example:} The group developed a devious scheme to obtain confidential information.

\textbf{Pattern:} devious tactics or methods

\textbf{Use:} Use this for dishonest or manipulative ways of achieving a goal.

\textbf{Example:} The investigation revealed that the firm had used devious tactics to eliminate competitors.

\textbf{Pattern:} a devious person

\textbf{Use:} Use this for someone who regularly uses dishonest or indirect methods.

\textbf{Example:} He was a devious businessman who rarely revealed his true intentions.

\textbf{Pattern:} a devious motive

\textbf{Use:} Use this when someone's hidden reason for doing something is dishonest or manipulative.

\textbf{Example:} She wondered whether there was a devious motive behind the generous offer.

\subsubsection{Common wrong usage}
\paragraph{Correction 1}
\textbf{Wrong:} He is devious because he likes solving difficult puzzles.

\textbf{Correct:} He is clever because he likes solving difficult puzzles.

\textbf{Why:} Devious is negative and suggests dishonesty or manipulation, not simply intelligence.

\paragraph{Correction 2}
\textbf{Wrong:} The company used deviously tactics.

\textbf{Correct:} The company used devious tactics.

\textbf{Why:} Use the adjective devious before the noun tactics.

\paragraph{Correction 3}
\textbf{Wrong:} She made a devious because she wanted to win.

\textbf{Correct:} She made a devious plan because she wanted to win.

\textbf{Why:} Devious is an adjective, so it normally needs a noun in this position.

\paragraph{Correction 4}
\textbf{Wrong:} His devious behavior was honest and completely transparent.

\textbf{Correct:} His behavior was honest and completely transparent.

\textbf{Why:} Devious normally implies hidden, dishonest, or manipulative behavior, so it conflicts with honest and transparent.

\subsection{Meaning 2: Indirect or winding}
\textbf{Part of speech:} adjective

\textbf{IPA:} /ˈdiːviəs/

\textbf{Pronunciation phrase:} devious

\textbf{Local audio file:} audios/devious.mp3

\textbf{Register:} formal and less common than the dishonest meaning; mainly used for routes, paths, and courses

\textbf{Definition:} indirect and complicated rather than following the simplest or most direct route

\textbf{In simpler words:} indirect and winding

\textbf{Typical pattern:} devious + route, path, course, or way

\subsubsection{Natural examples}
\begin{enumerate}
  \item Because the main road was closed, we had to take a devious route through several small villages.
  \item The river follows a devious course through the mountains.
  \item The hikers reached the valley by a long and devious path.
  \item The old road takes a devious route around the steepest part of the hill.
  \item Explorers followed a devious trail through the dense forest.
  \item The stream takes a devious course before eventually reaching the lake.
\end{enumerate}

\subsubsection{Nuance and implication}
\textbf{Central nuance:} This sense describes physical movement that does not follow a direct line.

\textbf{Usual implication:} The route may contain many turns or require extra distance.

\textbf{Compared with direct:} A direct route follows the shortest or simplest path. A devious route turns away from that direct path.

\subsubsection{Grammar notes}
\textbf{How to use it:} Use it before nouns such as route, path, course, or way.

\textbf{What can follow it:} It most naturally describes a route or physical course rather than a person.

\textbf{Position in a sentence:} This meaning is less common in everyday English than the dishonest meaning.

\subsubsection{Useful patterns}
\textbf{Pattern:} a devious route

\textbf{Use:} Use this for a route that reaches a destination in an indirect way.

\textbf{Example:} The travelers took a devious route around the flooded area.

\textbf{Pattern:} a devious path

\textbf{Use:} Use this for a path with many turns or changes in direction.

\textbf{Example:} A devious path led through the forest to the old village.

\textbf{Pattern:} a devious course

\textbf{Use:} Use this for a river, road, or other line of movement that follows an indirect direction.

\textbf{Example:} The river follows a devious course across the plain.

\subsubsection{Common wrong usage}
\paragraph{Correction 1}
\textbf{Wrong:} We took a devious route because it was the shortest and most direct road.

\textbf{Correct:} We took a direct route because it was the shortest road.

\textbf{Why:} A devious route is indirect or winding, not the shortest and most direct route.

\paragraph{Correction 2}
\textbf{Wrong:} The road moved devious through the mountains.

\textbf{Correct:} The road followed a devious route through the mountains.

\textbf{Why:} Use the adjective devious before a noun such as route or course.

\section{Word family}
\sectionrule
\subsection{deviously}
\textbf{Part of speech:} adverb

\textbf{IPA:} /ˈdiːviəsli/

\textbf{Definition:} in a dishonest, secretive, or manipulative way

\textbf{Relationship to the headword:} This adverb describes an action performed in a devious way.

\textbf{Grammar pattern:} deviously + manipulate, conceal, obtain, or plan

\textbf{Usage note:} It is less common than devious.

\subsubsection{Examples}
\begin{enumerate}
  \item He deviously concealed important information from the other negotiators.
  \item The company was accused of deviously manipulating its pricing system.
\end{enumerate}

\subsection{deviousness}
\textbf{Part of speech:} uncountable noun

\textbf{IPA:} /ˈdiːviəsnəs/

\textbf{Definition:} the quality of being dishonest, secretive, or cleverly manipulative

\textbf{Relationship to the headword:} This noun names the dishonest or manipulative quality described by devious.

\textbf{Grammar pattern:} the deviousness of + person, method, plan, or scheme

\textbf{Usage note:} It is considerably less common than devious.

\subsubsection{Examples}
\begin{enumerate}
  \item The deviousness of the scheme surprised investigators.
  \item His reputation for deviousness made colleagues reluctant to trust him.
\end{enumerate}

\section{Common collocations}
\sectionrule
\subsection{devious tactics}
\textbf{Applies to:} Dishonest and cleverly indirect

\textbf{Meaning:} dishonest or manipulative methods used to achieve a goal

\textbf{Example:} The company was accused of using devious tactics against its competitors.

\subsection{devious scheme}
\textbf{Applies to:} Dishonest and cleverly indirect

\textbf{Meaning:} a complicated and dishonest plan

\textbf{Example:} Investigators uncovered a devious scheme involving several false companies.

\subsection{devious plan}
\textbf{Applies to:} Dishonest and cleverly indirect

\textbf{Meaning:} a secretive or dishonest plan designed to achieve an advantage

\textbf{Example:} The character develops a devious plan to gain control of the organization.

\subsection{devious methods}
\textbf{Applies to:} Dishonest and cleverly indirect

\textbf{Meaning:} dishonest or indirect ways of accomplishing something

\textbf{Example:} The organization used devious methods to obtain confidential information.

\subsection{devious strategy}
\textbf{Applies to:} Dishonest and cleverly indirect

\textbf{Meaning:} a clever but dishonest or manipulative strategy

\textbf{Example:} The negotiator adopted a devious strategy intended to confuse the opposing side.

\subsection{devious motive}
\textbf{Applies to:} Dishonest and cleverly indirect

\textbf{Meaning:} a hidden and dishonest reason for acting

\textbf{Example:} Investigators suspected that he had a devious motive for offering financial assistance.

\subsection{devious route}
\textbf{Applies to:} Indirect or winding

\textbf{Meaning:} an indirect route with many changes in direction

\textbf{Example:} The road takes a devious route through the mountains.

\subsection{devious path}
\textbf{Applies to:} Indirect or winding

\textbf{Meaning:} a winding or indirect path

\textbf{Example:} The hikers followed a devious path through the forest.

\subsection{devious course}
\textbf{Applies to:} Indirect or winding

\textbf{Meaning:} an indirect or winding line of movement

\textbf{Example:} The river follows a devious course toward the coast.

\section{Synonyms and antonyms}
\sectionrule
\subsection{Close synonyms}
\subsubsection{sly}
\textbf{Part of speech:} adjective

\textbf{Applies to:} Dishonest and cleverly indirect

\textbf{Shared meaning:} Both can describe clever behavior that is not completely open or honest.

\textbf{Difference:} Sly often describes clever behavior that is secretive or slightly dishonest and can sometimes sound playful. Devious is more strongly negative and suggests deliberate manipulation or dishonesty.

\textbf{Example:} He gave a sly smile after secretly changing the plan.

\subsubsection{underhanded}
\textbf{Part of speech:} adjective

\textbf{Applies to:} Dishonest and cleverly indirect

\textbf{Shared meaning:} Both describe dishonest methods that are not open or straightforward.

\textbf{Difference:} Underhanded strongly emphasizes unfair or dishonest behavior carried out secretly. Devious also emphasizes cleverness and indirectness.

\textbf{Example:} The company was accused of using underhanded methods to eliminate competition.

\subsubsection{cunning}
\textbf{Part of speech:} adjective

\textbf{Applies to:} Dishonest and cleverly indirect

\textbf{Shared meaning:} Both can describe clever strategies that are not obvious.

\textbf{Difference:} Cunning emphasizes skill in achieving goals through clever planning and may be positive or negative depending on context. Devious is normally negative and emphasizes dishonesty or manipulation.

\textbf{Example:} The animal uses a cunning strategy to obtain food.

\subsubsection{winding}
\textbf{Part of speech:} adjective

\textbf{Applies to:} Indirect or winding

\textbf{Shared meaning:} Both can describe a route that is not straight.

\textbf{Difference:} Winding simply describes something with many bends or turns. Devious is more formal and emphasizes that the route is indirect.

\textbf{Example:} A winding road climbs through the mountains.

\subsection{Direct or useful antonyms}
\subsubsection{straightforward}
\textbf{Part of speech:} adjective

\textbf{Applies to:} Dishonest and cleverly indirect

\textbf{Opposition:} Straightforward behavior is open, honest, and direct, while devious behavior relies on hidden or manipulative methods.

\textbf{Example:} She gave a straightforward explanation of what had happened.

\subsubsection{direct}
\textbf{Part of speech:} adjective

\textbf{Applies to:} Indirect or winding

\textbf{Opposition:} A direct route follows the simplest path toward a destination, while a devious route is indirect and winding.

\textbf{Example:} We took the most direct route to the village.

\section{Words learners may confuse}
\sectionrule
\subsection{deviant}
\textbf{Part of speech:} adjective

\textbf{Why learners confuse them:} The words look similar and both come from roots connected with turning away from a normal path.

\textbf{Key difference:} Devious means dishonest and manipulative or, less commonly, indirect and winding. Deviant describes behavior or a person that differs significantly from accepted social norms.

\textbf{devious:} The executive used devious tactics to hide the company's losses.

\textbf{deviant:} The study examined behavior considered deviant by the surrounding society.

\subsection{deviate}
\textbf{Part of speech:} adjective versus verb

\textbf{Why learners confuse them:} The words share a similar form and both contain the idea of moving away from a direct or expected path.

\textbf{Key difference:} Devious is an adjective describing dishonest behavior or an indirect route. Deviate is a verb meaning move away from an expected path, standard, or plan.

\textbf{devious:} They developed a devious plan to avoid the regulations.

\textbf{deviate:} The aircraft had to deviate from its original route because of bad weather.

\section{Etymology}
\sectionrule
\textbf{Origin:} Devious comes from Latin devius, meaning out of the way or away from the road.

\textbf{Development:} The original idea of leaving the direct path developed into the figurative meaning of using indirect, secretive, or dishonest methods.

\textbf{Memory link:} A devious person does not take the straight path toward a goal; they use indirect and often dishonest methods.

\section{Practice exercises}
\sectionrule
\subsection{Word family choice}
Choose the best member of the word family for each blank.

\begin{enumerate}
  \item The \_\_\_\_\_ of the scheme made it difficult for investigators to discover how it worked.
  \item The company \_\_\_\_\_ concealed several important fees from customers.
  \item Investigators uncovered a \_\_\_\_\_ plan to hide the money.
\end{enumerate}

\subsection{Correct or incorrect?}
Decide whether each sentence is correct. Correct the inaccurate ones.

\begin{enumerate}
  \item The company used devious tactics to hide information from customers.
  \item She is devious because she is very intelligent and solves problems quickly.
  \item The road follows a devious course through the mountains.
  \item The company used deviously methods to avoid the rule.
  \item His devious behavior was completely honest and transparent.
\end{enumerate}

\subsection{Collocation practice}
Complete each sentence with a natural collocation.

\begin{enumerate}
  \item Investigators uncovered a devious \_\_\_\_\_ designed to hide the company's profits.
  \item The politician was accused of using devious \_\_\_\_\_ to mislead voters.
  \item Because the bridge was closed, we had to take a devious \_\_\_\_\_ through the hills.
  \item She suspected that he had a devious \_\_\_\_\_ for offering to help.
\end{enumerate}

\subsection{Complete answer key}
\subsubsection{Word family choice}
\begin{enumerate}
  \item \textbf{deviousness.} The deviousness of the scheme made it difficult for investigators to discover how it worked. \emph{Use the noun deviousness for the quality of being dishonest or manipulative.}
  \item \textbf{deviously.} The company deviously concealed several important fees from customers. \emph{Use the adverb deviously to describe how the company concealed the fees.}
  \item \textbf{devious.} Investigators uncovered a devious plan to hide the money. \emph{Use the adjective devious before plan.}
\end{enumerate}

\subsubsection{Correct or incorrect?}
\begin{enumerate}
  \item \textbf{Correct} \emph{Devious tactics naturally describes dishonest or manipulative methods.}
  \item \textbf{Incorrect}. She is clever because she is very intelligent and solves problems quickly. \emph{Devious has a negative meaning and suggests dishonesty or manipulation, not merely intelligence.}
  \item \textbf{Correct} \emph{Devious can describe an indirect or winding physical route.}
  \item \textbf{Incorrect}. The company used devious methods to avoid the rule. \emph{Use the adjective devious before the noun methods.}
  \item \textbf{Incorrect}. His behavior was completely honest and transparent. \emph{Devious behavior is normally dishonest, secretive, or manipulative.}
\end{enumerate}

\subsubsection{Collocation practice}
\begin{enumerate}
  \item \textbf{scheme.} Investigators uncovered a devious scheme designed to hide the company's profits. \emph{A devious scheme is a dishonest and carefully planned arrangement.}
  \item \textbf{tactics.} The politician was accused of using devious tactics to mislead voters. \emph{Devious tactics are dishonest or manipulative methods.}
  \item \textbf{route.} Because the bridge was closed, we had to take a devious route through the hills. \emph{A devious route is indirect and winding.}
  \item \textbf{motive.} She suspected that he had a devious motive for offering to help. \emph{A devious motive is a hidden or dishonest reason for doing something.}
\end{enumerate}

\vfill
\begin{center}
\small\color{Muted} Created and compiled by \textbf{Amir Kooshky}\\
\href{https://t.me/KooshkyTOEFL}{Telegram Channel}\quad\textbullet\quad
\href{https://www.instagram.com/kooshkytoefl}{Instagram}\quad\textbullet\quad
\href{https://t.me/KooshkyTOEFL\_pv}{Personal Telegram}
\end{center}
\end{document}
